\documentclass[11pt]{amsart}
\usepackage[T1]{fontenc}
\usepackage{lmodern,microtype,amsmath,amssymb,amsthm,mathtools}
\usepackage[margin=1.1in]{geometry}
\usepackage[hidelinks]{hyperref}
\usepackage{xurl}

\newtheorem{theorem}{Theorem}[section]
\newtheorem{corollary}[theorem]{Corollary}
\newtheorem{lemma}[theorem]{Lemma}
\newtheorem{proposition}[theorem]{Proposition}
\theoremstyle{remark}
\newtheorem{remark}[theorem]{Remark}
\numberwithin{equation}{section}

\newcommand{\T}{\mathbb{T}}
\newcommand{\Z}{\mathbb{Z}}
\newcommand{\chiB}{\chi_{\mathrm B}}
\newcommand{\chiM}{\chi_{\mu}}

\hypersetup{pdftitle={Borel Three-Colorings of Unit-Distance Graphs on Two Concentric Circles},pdfauthor={Yitzchak Shmalo}}

\begin{document}

\title[Three-colorings on two concentric circles]{Borel Three-Colorings of Unit-Distance Graphs on Two Concentric Circles}
\author{Yitzchak Shmalo}
\address{Einstein Institute of Mathematics, The Hebrew University of Jerusalem, Jerusalem, Israel}
\email{yitzchak.shmalo@gmail.com}
\date{}
\subjclass[2020]{05C15, 03E15, 52C10, 11R18}
\keywords{unit-distance graph, Borel chromatic number, measurable chromatic number, rose window graph, irrational rotation, cyclotomic units}

\begin{abstract}
Take two concentric circles on which a chord of length one subtends $q$
turns and $p$ turns, where $q<p$ are rational, and join two points of their union when
they are at distance one. Suppose that $p\equiv\pm2q\pmod 1$ and that
the angle $\pi g$ at the centre between the endpoints of a unit segment
joining the two circles is an irrational multiple of $\pi$. We show that
the chromatic number, the Lebesgue measurable chromatic number and the
Borel chromatic number of this graph all equal $3$. The coloring is
constant on finitely many arcs and comes from a $3$-coloring
of one of Wilson's rose window graphs. Under the same hypotheses, the
coloring $x\mapsto\lfloor3mx\rfloor\bmod3$ of the turn $x$ on both
circles is proper for some $m$ exactly when $3$ divides the denominator
of $q$. The irrationality hypothesis holds whenever this denominator is
twice an odd prime power or a power of $2$ at least $4$, giving
infinitely many examples.
\end{abstract}

\maketitle

\section{Introduction}

Let $X$ be a set in the plane and let $G(X)$ be the graph on $X$ in
which two points are adjacent when they are at Euclidean distance one.
For $X=\mathbb R^2$ this is the graph of the Hadwiger--Nelson
problem~\cite{Soi09}. For a set $X$ carrying a Borel structure and a
measure one can also ask for proper colorings that are Borel or
measurable. The resulting chromatic numbers were introduced in a general
setting by Kechris, Solecki and Todorcevic~\cite{KST99}, and they can be
larger than the ordinary chromatic number. The basic example, recalled
by Conley, Marks and Tucker-Drob~\cite{CMT16}, is the graph on the
circle that joins each point to its image under an irrational rotation.
It is a union of bi-infinite paths, so it has chromatic number $2$, but
it has no measurable proper $2$-coloring.

We study a simple planar instance in which irrational rotation
appears naturally. For $0<t\le 1/2$ put $R(t)=1/(2\sin\pi t)$, so that a
unit chord of the circle of radius $R(t)$ subtends the angle $2\pi t$.
Given $0<q<p\le1/2$ let
\[
r=R(q),\qquad s=R(p),
\]
and let $U(q,p)$ be the unit-distance graph on the union of the circle
of radius $r$ and the circle of radius $s$ about the origin. Since
$q<p$ we have $r>s$. Two points on different circles are joined by an
edge, which we call a \emph{cross edge}, when their angular difference
is $\pm\pi g$, where $g\in[0,1]$ is defined by
\begin{equation}\label{eq:cross}
r^2+s^2-2rs\cos(\pi g)=1.
\end{equation}
We call the pair \emph{compatible} when this equation has a solution.
Measuring angles in turns, a point at turn $x$ on the outer circle is
adjacent to $x\pm q$ on the same circle and to $x\pm g/2$ on the inner
circle, and a point at turn $y$ on the inner circle is adjacent to
$y\pm p$ on the inner circle. When $q$ and $p$ are rational and $g$ is
irrational, the rotations by $q$ and by $p$ have finite order and the
cross edges carry an irrational rotation.

We write $\chi$, $\chiM$ and $\chiB$ for the chromatic number, the
Lebesgue measurable chromatic number and the Borel chromatic number of
$U(q,p)$. Here the measure is arc length on each circle, and a
coloring is required to be proper at every point; Remark~\ref{rem:null}
shows that the lower bound does not change if properness is only
required off a null set.

The general theorems give weaker bounds for these graphs. Every vertex of $U(q,p)$ has degree at most $4$ and there is no
$K_5$, so the measurable Brooks theorem of Conley, Marks and
Tucker-Drob~\cite[Theorem~1.2]{CMT16} gives a measurable proper
$4$-coloring in their sense. Under the hypotheses of
Theorem~\ref{thm:main} below, every connected component of $U(q,p)$ is
two-ended. With $q=a/n$ in lowest terms, the turns of the points of a
component have the form $\theta+i/n+jg/2$ with $i,j\in\Z$ and $\theta$
fixed. Since $g$ is irrational, the point determines $j$, and each value
of $j$ is taken by at most $n$ points. The rotation by $g$ maps each
component to itself, since $A(x)$ and $A(x+g)$ are both joined to
$B(x+g/2)$, and likewise on the inner circle. It raises $j$ by $2$, so it
acts freely, with at most $2n$ orbits on a component. A connected locally
finite graph with a free action of $\Z$ with finitely many orbits is
two-ended. For such graphs Weilacher's
theorem~\cite[Theorem~1.3]{Wei20} gives $\chiB\le2\chi-1$. Greb\'ik and
Rozho\v{n}~\cite{GR21} classify local labelling problems with finite
input labels on aperiodic Borel actions of $\Z$. After each point is
grouped with its rotations by multiples of $1/n$, a proper $3$-coloring
of $U(q,p)$ becomes such a problem over the rotation by $g/2$ that
exchanges the two circles. Their Theorem~5.9 gives a Borel solution
whenever the problem is not mixing in the sense
of~\cite[Definition~5.8]{GR21}. We do not decide whether this problem is
mixing. Theorem~\ref{thm:main} below gives $\chiM=\chiB=3$ either way,
by a coloring that is constant on finitely many arcs and given by the
same construction for every $n$.

\begin{theorem}\label{thm:main}
Let $0<q<p\le1/2$ be rational with $p\equiv\pm2q\pmod1$, suppose the
pair is compatible, and suppose $g$ is irrational. Then
\[
\chi(U(q,p))=\chiM(U(q,p))=\chiB(U(q,p))=3,
\]
and $U(q,p)$ has a proper $3$-coloring that is constant on each of
finitely many half-open arcs.
\end{theorem}

With $\chi=3$, the bounds above are $4$ for $\chiM$ and $5$ for $\chiB$.
The coloring of Theorem~\ref{thm:main} is explicit, \eqref{eq:color}
below. Write $q=a/n$ in lowest terms. The coloring comes from a proper
$3$-coloring of the finite graph $T_n$ on the vertices $A_j,B_j$
($j\in\Z/n\Z$) with edges $A_jA_{j+1}$, $B_jB_{j+2}$, $A_jB_j$ and
$A_jB_{j+1}$. This is Wilson's rose window graph $R_n(1,2)$~\cite{Wil08}
with its indices reversed, and it is $3$-colorable for every $n\ge4$
except $n=5$ (Lemma~\ref{lem:palette}). Both circles are cut into
$L=nk$ equal arcs, with $k\equiv a^{-1}\pmod n$ chosen so that
\[
Lq\equiv1\pmod n,\qquad Lp\equiv\pm2\pmod n,\qquad
\frac{Lg}{2}=vn+\varepsilon\quad(v\in\Z,\ |\varepsilon|<\tfrac12).
\]
The rotations by $q$ and $p$ then move the index of an arc by $1$ and by
$\pm2$ modulo $n$. Once the arcs of the inner circle are offset by half
an arc, the rotations by $\pm g/2$ take the arc of index $j$ on the outer
circle into the arcs of indices $j$ and $j+1$ modulo $n$ on the inner
circle. So every edge of $U(q,p)$ lands on an edge of $T_n$. The lower
bound comes from a $5$-cycle. The condition $p\equiv\pm2q$ forces
$q<1/3$, and compatibility holds exactly for $q_0\le q<1/3$, where
$q_0\approx0.0779$ is defined in Lemma~\ref{lem:compat}.

Theorem~\ref{thm:main} needs $g$ to be irrational, and for small
denominators this sometimes fails: among the $284$ compatible pairs with
denominator of $q$ at most $60$, the number $g$ is rational exactly for
$q=1/5,1/7,2/7,1/9,2/9$ and $1/12$ (Section~\ref{sec:comp}). The next
result gives an infinite family where irrationality is guaranteed.

\begin{theorem}\label{thm:family}
Let $q=a/n$ in lowest terms with $q_0\le q<1/3$, and let $p$ be the
element of $(q,1/2]$ congruent to $\pm2q$. If $n=2\ell^e$ for an odd
prime $\ell$ and $e\ge1$, or $n=2^e$ with $e\ge2$, then $g$ is
irrational. Hence Theorem~\ref{thm:main} applies to infinitely many
pairs, among them $q=1/10$, $p=1/5$.
\end{theorem}

By Lemma~\ref{lem:compat}, $2\cos(\pi g)=\xi^3-3\xi+1/\xi$ with
$\xi=2\cos(\pi q)$, and $\xi$ is not a unit exactly for these $n$, so
that $2\cos(\pi g)$ is not an algebraic integer. For $q=1/10$, $p=1/5$
one may take $110$ arcs on each circle, and Corollary~\ref{cor:tenth}
gives the coloring explicitly.

A natural first attempt is a \emph{sector coloring}: cut the plane into
$3m$ equal sectors about the origin and color them cyclically, so that
the point at turn $x$ on either circle gets the color
$\lfloor3mx\rfloor\bmod3$. This is a radial coloring in the sense of
Alm and Manske~\cite{AM14}: there are finitely many rays from the
origin, each open sector between consecutive rays has one color, and
the rays themselves may be colored arbitrarily. Alm and Manske study radial colorings of the
unit-distance graphs of annuli of outer radius less than $1$. Since
$r<1$ exactly when $q>1/6$, for $1/6<q<1/3$ the graph $U(q,p)$ lies in
such an annulus. Chybowska-Sok\'o\l, Junosza-Szaniawski and
W\k{e}sek~\cite[Proposition~4.1]{CJW22} use radial colorings to bound
the chromatic number of the graph on an annulus that joins points at
distances in an interval $[1,b]$. Proposition~\ref{prop:common} shows
that under the hypotheses of Theorem~\ref{thm:main} a sector coloring
of $U(q,p)$ exists exactly when $3$ divides $n$, and
Remark~\ref{rem:radial} shows that for $3\nmid n$ no proper
$3$-coloring of $U(q,p)$ uses the same function of the turn on both
circles. So sector
colorings miss, for example, $q=1/10$, and Theorem~\ref{thm:family}
gives infinitely many pairs that they miss.

The facts about cyclotomic units used in Section~\ref{sec:family} are
standard~\cite[Chapter~2]{Was97}. For the pairs with $3\nmid n$, among
them $q=1/10$ and infinitely many pairs of Theorem~\ref{thm:family},
neither a sector coloring
nor the general bounds of~\cite{CMT16} and~\cite{Wei20} give the value
$3$ for $\chiM$ and $\chiB$. The theorem of Greb\'ik and
Rozho\v{n}~\cite[Theorem~5.9]{GR21} gives a Borel proper $3$-coloring
for these pairs once the coloring problem is shown not to be mixing, and
with the lower bound from a $5$-cycle this gives $\chiM=\chiB=3$. We do
not check the mixing condition directly; Theorem~\ref{thm:main} gives
an explicit coloring.

Section~\ref{sec:model} computes the cross angle $g$ and the range of
compatible pairs (Lemma~\ref{lem:compat}). Section~\ref{sec:palette}
gives a proper $3$-coloring of $T_n$ for $n\ge4$, $n\ne5$
(Lemma~\ref{lem:palette}), and Section~\ref{sec:proof} pulls it back to
$U(q,p)$, which proves Theorem~\ref{thm:main}.
Section~\ref{sec:common} shows that a sector coloring exists exactly
when $3\mid n$ (Proposition~\ref{prop:common}).
Section~\ref{sec:family} proves that $g$ is irrational for the family
of Theorem~\ref{thm:family} and writes out the coloring for $q=1/10$.
Section~\ref{sec:comp} lists the rational values of $g$ for $n\le60$
and the arc-by-arc check of \eqref{eq:color}, and
Section~\ref{sec:questions} asks whether these are all the rational
values, what happens when $p\not\equiv\pm2q$, and how large $k$ must
be.

\section{The model}\label{sec:model}

We use the turn coordinate $x\in\T=\mathbb R/\Z$ on each circle. A point
of $U(q,p)$ is written $A(x)$ if it lies on the outer circle (radius
$r$) and $B(x)$ if it lies on the inner circle (radius $s$). Two points
on the circle of radius $R(t)$ at turns $x,y$ are at distance
$2R(t)\sin(\pi\|x-y\|)$, where $\|\cdot\|$ is the distance to $\Z$; this
equals $1$ exactly when $x-y\equiv\pm t$. A point $A(x)$ and a point
$B(y)$ are at distance one exactly when $\cos(2\pi(x-y))=\cos(\pi g)$,
that is, $y-x\equiv\pm g/2$. So the edges of $U(q,p)$ are
\begin{equation}\label{eq:edges}
A(x)A(x\pm q),\qquad B(x)B(x\pm p),\qquad A(x)B(x\pm g/2).
\end{equation}
We call the first two kinds \emph{chord edges}; the third kind are the
cross edges. When $p=1/2$ the two chord neighbours of a point on the
inner circle coincide.

\begin{lemma}\label{lem:compat}
Let $0<q<p\le1/2$ with $p\equiv\pm2q\pmod1$. Then $q<1/3$, and
$p=2q$ if $q\le1/4$, $p=1-2q$ if $1/4<q<1/3$. With $\xi =2\cos(\pi q)\in(1,2)$,
\begin{equation}\label{eq:g}
s=\frac{r}{\xi},\qquad \frac{r^2+s^2-1}{2rs}=\cos(3\pi q)+\frac{1}{4\cos(\pi q)},
\end{equation}
so a solution of \eqref{eq:cross} satisfies
\[
2\cos(\pi g)=\xi^3-3\xi +\frac1\xi.
\]
The pair is compatible exactly when $\xi\le \xi_0$, where $\xi_0$ is the
unique root in $(1,2)$ of $X^4-3X^2-2X+1$. Equivalently $q\ge q_0$ with
$q_0=\arccos(\xi_0/2)/\pi$. Rounded to five decimals, $\xi_0\approx1.94039$
and $q_0\approx0.07791$.
\end{lemma}

\begin{proof}
If $q\ge1/3$ then $2q\in[2/3,1)$, so the representatives of $\pm2q$ in
$[0,1/2]$ are $1-2q\le1/3\le q$, and no $p\in(q,1/2]$ exists. For
$q<1/3$ the representative of $\pm2q$ in $(0,1/2]$ is $2q$ or $1-2q$
as stated, and it exceeds $q$. In both cases $\sin(\pi p)=\sin(2\pi q)$,
so $s=1/(2\sin2\pi q)=r/(2\cos\pi q)=r/\xi$. Put $c=\cos\pi q$. Using
$1/r^2=4\sin^2\pi q=4(1-c^2)$,
\[
\frac{r^2+s^2-1}{2rs}=\frac{r^2(1+1/(4c^2))-1}{r^2/c}
=c+\frac{1}{4c}-4c(1-c^2)=4c^3-3c+\frac1{4c},
\]
which is \eqref{eq:g}, since $4c^3-3c=\cos 3\pi q$. Put
$f(\xi)=\xi^3-3\xi +1/\xi$. For $\xi\in(1,2)$,
\[
\xi\,(f(\xi)+2)=(\xi^2-1)^2+\xi(2-\xi)>0,\qquad \xi\,(f(\xi)-2)=\xi^4-3\xi^2-2\xi +1=:h(\xi).
\]
So $f>-2$ on $(1,2)$, and $f(\xi)\le2$ exactly when $h(\xi)\le0$. Now
$h(1)=-3$, $h(2)=1$ and $h''(\xi)=12\xi^2-6>0$ on $(1,2)$, so $h$ has a
unique root $\xi_0$ in $(1,2)$ and is negative to its left. Equation
\eqref{eq:cross} has a solution $g\in[0,1]$ exactly when
$|f(\xi)|\le2$, which is $\xi\le \xi_0$.
\end{proof}

When $g$ is irrational we have $0<g<1$, and the two cross neighbours of
a point are distinct. Each vertex then has degree $4$, or degree $3$ on
the inner circle when $p=1/2$.

\section{A three-coloring of the rose window graph}\label{sec:palette}

For $n\ge4$ let $T_n$ be the graph with vertices $A_j,B_j$,
$j\in\Z/n\Z$, and edges
\begin{equation}\label{eq:Tn}
A_jA_{j+1},\qquad B_jB_{j+2},\qquad A_jB_j,\qquad A_jB_{j+1},
\end{equation}
parallel edges being identified (for $n=4$ the edges $B_jB_{j+2}$ and
$B_{j+2}B_j$ coincide). In Wilson's notation~\cite{Wil08} the rose
window graph $R_n(u,w)$ has rim edges $A_iA_{i+1}$, spokes $A_iB_i$
and $A_{i+u}B_i$, and hub edges $B_iB_{i+w}$; the substitution
$j\mapsto-j$ carries $T_n$ to $R_n(1,2)$. Colors are elements of
$\Z/3\Z$.

\begin{lemma}\label{lem:palette}
Let $n\ge4$, $n\ne5$. Put $m_-=3,2,4$ according as $n\equiv0,1,2\pmod3$,
and $m_+=n-m_-$. Let $d_j=+1$ for $0\le j<m_+$ and $d_j=-1$ for
$m_+\le j<n$, indices modulo $n$, and define
\begin{equation}\label{eq:palette}
\alpha_0=0,\qquad \alpha_{j+1}=\alpha_j+d_j,\qquad \beta_j=-\alpha_{j-1}-\alpha_j\qquad(\text{in }\Z/3\Z).
\end{equation}
Then $\alpha_n=\alpha_0$, and coloring $A_j$ by $\alpha_j$ and $B_j$ by $\beta_j$ is a proper
$3$-coloring of $T_n$.
\end{lemma}

\begin{proof}
The sum of the $d_j$ is $m_+-m_-=n-2m_-$, which is divisible by $3$ in
each of the three cases, so the sequence $\alpha_j$ closes up. Both
$m_+$ and $m_-$ are at least $2$: this holds for $n=4,6,7,8$ by
inspection ($m_+=2,3,5,4$) and $m_+$ only grows with $n$ in each residue
class. For $n=5$ we would get $m_+=1$, which is why it is excluded.

Consecutive values $\alpha_{j-1},\alpha_j$ differ by $\pm1$, and since the three
elements of $\Z/3\Z$ sum to $0$, the value $\beta_j$ is the third color,
different from both. Hence $\beta_j\ne \alpha_j$ and $\beta_{j+1}\ne \alpha_j$, so both
kinds of spokes are proper, and the rim $A_jA_{j+1}$ is proper since
$d_j\ne0$. For the hub, $\alpha_{j-1}=\alpha_j-d_{j-1}$ gives
$\beta_j=-2\alpha_j+d_{j-1}=\alpha_j+d_{j-1}$, and therefore
\begin{equation}\label{eq:hub}
\beta_{j+2}-\beta_j=d_j+2d_{j+1}-d_{j-1}=d_j-d_{j+1}-d_{j-1}.
\end{equation}
If $d_{j-1}=-d_{j+1}$ the right side is $d_j\ne0$. If
$d_{j-1}=d_{j+1}=\sigma$ it is $d_j-2\sigma=d_j+\sigma$, which
vanishes only when $d_j=-\sigma$, that is, when $d_j$ is a run of
length one in the cyclic sequence $(d_j)$. Both runs have length at least
two, so $\beta_{j+2}\ne \beta_j$ for every $j$.
\end{proof}

\begin{remark}\label{rem:five}
The exclusion of $n=5$ is necessary: $T_5$ has no proper $3$-coloring.
In such a coloring the increments $d_j=\alpha_{j+1}-\alpha_j\in\{\pm1\}$ sum to $0$
modulo $3$, which for five signs forces one sign to occur exactly once.
The triangles $A_{j-1}A_jB_j$ force $\beta_j=-\alpha_{j-1}-\alpha_j$, and
\eqref{eq:hub} then makes the hub edge at the isolated sign
monochromatic. The case does not arise in Theorem~\ref{thm:main}: by
Lemma~\ref{lem:compat}, $n=5$ forces $q=1/5$ and $p=2/5$, and then
$r^2+s^2=1$ by the identity
$\sin^2(\pi/5)+\sin^2(2\pi/5)=5/4$ together with
$4\sin^2(\pi/5)\sin^2(2\pi/5)=5/4$, so $\cos(\pi g)=0$ and $g=1/2$ is
rational.
\end{remark}

\section{Proof of Theorem~\ref{thm:main}}\label{sec:proof}

Let $q=a/n$ in lowest terms. By Lemma~\ref{lem:compat}, $q<1/3$, so
$n\ge4$, and $n\ne5$ by Remark~\ref{rem:five} because $g$ is
irrational. Let $(\alpha_j)$, $(\beta_j)$ be the sequences of
Lemma~\ref{lem:palette}.

\emph{Choice of the number of arcs.} Let $k_0\in\{1,\dots,n-1\}$ with
$k_0a\equiv1\pmod n$. Since $ng/2$ is irrational, the numbers
$(k_0+nt)g/2$, $t=0,1,2,\dots$, are dense modulo $1$ (the forward orbit
of an irrational rotation is dense). So there is $k\equiv k_0\pmod n$,
$k\ge1$, with $\|kg/2\|<1/(2n)$. Put $L=nk$. Then
\begin{equation}\label{eq:L}
Lq=ka\equiv1\pmod n,\qquad Lp\equiv\pm2\pmod n,\qquad
\frac{Lg}{2}=vn+\varepsilon\ \text{ with } v\in\Z,\ |\varepsilon|<\tfrac12.
\end{equation}
The second relation holds because $p\equiv\pm2q\pmod1$ and $n\mid L$.

\emph{The coloring.} Define
\begin{equation}\label{eq:color}
C(A(x))=\alpha_{\lfloor Lx\rfloor\bmod n},\qquad
C(B(x))=\beta_{\lfloor Lx+1/2\rfloor\bmod n}.
\end{equation}
These are well defined on $\T$ because replacing $x$ by $x+1$ adds $L$,
a multiple of $n$, to both indices. Each is constant on $L$ half-open
arcs of each circle, so $C$ is Borel and Lebesgue measurable.

\emph{Chord edges.} Since $Lq$ and $Lp$ are integers,
$\lfloor L(x+q)\rfloor=\lfloor Lx\rfloor+Lq$ and
$\lfloor L(x+p)+1/2\rfloor=\lfloor Lx+1/2\rfloor+Lp$. By \eqref{eq:L}
the edge $A(x)A(x+q)$ receives the colors of $A_j,A_{j+1}$ and the edge
$B(x)B(x+p)$ the colors of $B_i,B_{i\pm2}$, for suitable $j,i$. These
are edges of $T_n$, so their colors differ by Lemma~\ref{lem:palette}.

\emph{Cross edges.} Write $Lx=j+u$ with $j\in\Z$ and $0\le u<1$. Then
\begin{equation}\label{eq:crossindex}
\Bigl\lfloor L\bigl(x\pm \tfrac g2\bigr)+\tfrac12\Bigr\rfloor
=j\pm vn+\bigl\lfloor u+\tfrac12\pm\varepsilon\bigr\rfloor ,
\end{equation}
and $0\le u+\frac12\pm\varepsilon<2$, so the last term is $0$ or $1$.
Hence the cross edge from $A(x)$ receives the colors of $A_j$ and of
$B_j$ or $B_{j+1}$ (indices modulo $n$), which differ by
Lemma~\ref{lem:palette}. So $C$ is a proper $3$-coloring.

\emph{Lower bound.} The five points
\begin{equation}\label{eq:cycle}
A(x),\ A(x+q),\ A(x+2q),\ B(x+2q+g/2),\ B(x+g/2)
\end{equation}
form a closed walk of length five in $U(q,p)$: the first two steps are
chord steps by $q$ on the outer circle, the third and fifth are cross
steps by $g/2$, and the fourth is a chord step on the inner circle, by
$2q\equiv\pm p$. The three outer points
are distinct because $q$ has order $n\ge4$, and the two inner points
are distinct because $p\not\equiv0$. So \eqref{eq:cycle} is a $5$-cycle
and $\chi\ge3$. Since $\chi\le\chiM\le\chiB\le3$, all three equal $3$.
\qed

\begin{remark}\label{rem:null}
If a measurable coloring is only required to be proper outside a set of
measure zero, the lower bound still holds. The five maps sending $x$ to
the points of \eqref{eq:cycle} are rotations, so they preserve arc length.
Hence for almost every $x$ all five points avoid a given null set, and a
$2$-coloring that is proper off that set would be proper on the
$5$-cycle at such an $x$, which is impossible.
\end{remark}

\begin{remark}\label{rem:hyp}
The coloring \eqref{eq:color} is proper whenever $n\ne5$ and $k\ge1$
satisfies $ka\equiv1\pmod n$ and $\|kg/2\|\le1/(2n)$, since
\eqref{eq:crossindex} also covers $|\varepsilon|=\frac12$; irrationality
of $g$ serves only to exclude $n=5$ and to provide such a $k$. The shift by
half an arc on the inner circle is needed in \eqref{eq:color}: for
$q=1/10$, $k=11$, using $\lfloor Lx\rfloor$ on both circles produces
monochromatic cross edges, and so does $k=2$ (which violates
$ka\equiv1$) with the shift kept. Section~\ref{sec:comp} gives the
numbers of monochromatic pairs of arcs.
\end{remark}

\section{Sector colorings}\label{sec:common}

\begin{proposition}\label{prop:common}
Under the hypotheses of Theorem~\ref{thm:main}, with $q=a/n$ in lowest
terms, there is a positive integer $m$ such that
$x\mapsto\lfloor3mx\rfloor\bmod3$, used on both circles, is a proper
coloring of $U(q,p)$ if and only if $3$ divides $n$.
\end{proposition}

\begin{proof}
Let $D(x)=\lfloor3mx\rfloor\bmod3$ and let $\tau\in\mathbb R$. Write
$z$ for the representative of $3m\tau$ in $[0,3)$. For $3mx=i+u$ with
$i\in\Z$, $0\le u<1$, we get $D(x+\tau)-D(x)\equiv\lfloor u+z\rfloor\pmod3$.
If $z<1$, take $u=0$; if $z>2$, take $u=3-z\in(0,1)$; in both cases the
difference is $0$ modulo $3$. If $1\le z\le2$ the difference is $1$ or
$2$ for every $u$. Hence $D(x)\ne D(x+\tau)$ for all $x$ exactly when
$\{m\tau\}\in[1/3,2/3]$, where $\{\cdot\}$ is the fractional part. Note
that this set is invariant under $y\mapsto1-y$.

By \eqref{eq:edges}, $D$ is proper exactly when $\{m\tau\}\in[1/3,2/3]$
for $\tau=q$, $\tau=p$ and $\tau=g/2$. Put $y=\{mq\}$. Since
$p\equiv\pm2q$, the condition for $\tau=p$ is $\{2y\}\in[1/3,2/3]$. If
$y\in[1/3,1/2)$ then $\{2y\}=2y\in[2/3,1)$, so $y=1/3$. If
$y\in[1/2,2/3]$ then $\{2y\}=2y-1\in[0,1/3]$, so $y=2/3$. Thus the two
chord conditions hold exactly when $\{ma/n\}\in\{1/3,2/3\}$, which is
possible only if $3\mid n$.

Conversely, if $3\mid n$, pick $m_0\in\{1,\dots,n-1\}$ with $m_0a\equiv n/3\pmod n$. Every
$m=m_0+nt$ with $t\ge0$ satisfies $\{mq\}=1/3$, and since $ng/2$ is
irrational the numbers $mg/2$ are dense modulo $1$; choose $t$ with
$\{mg/2\}\in(1/3,2/3)$.
\end{proof}

\begin{remark}\label{rem:radial}
The condition $3\mid n$ is needed for every proper $3$-coloring of
$U(q,p)$ that colors a point by the same function of its turn on both
circles, not only for sector colorings. (The coloring of
Theorem~\ref{thm:main} uses different functions on the two circles.)
Let $p\equiv\pm2q$, and suppose $F:\T\to\Z/3\Z$ is such that coloring
both $A(x)$ and $B(x)$ by $F(x)$ is proper. The edges $A(x)A(x+q)$,
$A(x+q)A(x+2q)$ and $B(x)B(x+2q)$ show that $F(x)$, $F(x+q)$ and
$F(x+2q)$ are distinct for every $x$. The same holds with $x$ replaced
by $x+q$, so $F(x+3q)=F(x)$ for every $x$. If $3\nmid n$, then $3q$
generates the same subgroup of $\T$ as $q$, so $F(x+q)=F(x)$, and the
edge $A(x)A(x+q)$ is monochromatic. The argument also applies when $F$
is defined only off finitely many turns, as for a radial coloring
in the sense of~\cite{AM14}: for all but finitely many $x$ the turns
$x+i/n$, $i\in\Z$, avoid them.
\end{remark}

\section{Irrationality of \texorpdfstring{$g$}{g}}\label{sec:family}

\begin{proof}[Proof of Theorem~\ref{thm:family}]
By Lemma~\ref{lem:compat} the pair is compatible, and
$2\cos(\pi g)=\xi^3-3\xi+1/\xi$ with $\xi=2\cos(\pi a/n)$. If $g$ were rational,
$2\cos(\pi g)=\zeta+\zeta^{-1}$ for a root of unity $\zeta$ would be an
algebraic integer. Since $\xi$ is an algebraic integer, so would be $1/\xi$,
and $\xi$ would be a unit. We show that $\xi$ is not a unit for the stated $n$.

Let $\rho=e^{\pi i/n}$ and $\omega=\rho^{2a}=e^{2\pi ia/n}$, a primitive
$n$-th root of unity since $\gcd(a,n)=1$. Then $\xi=\rho^{-a}(1+\omega)$,
so $\xi$ is a unit exactly when $1+\omega=(1-\omega^2)/(1-\omega)$ is. We
use the standard facts~\cite[Chapter~2]{Was97}: if $\eta$ is a primitive
$d$-th root of unity with $d>1$, then the norm of $1-\eta$ from
$\mathbb Q(\eta)$ to $\mathbb Q$ is $\Phi_d(1)$, where $\Phi_d$ is the
$d$-th cyclotomic polynomial, and $\Phi_d(1)$ equals $\ell$ if
$d$ is a power of a prime $\ell$ and $1$ otherwise. We compute norms
from $K=\mathbb Q(\omega)$, of degree $\varphi(n)$, where $\varphi$ is
Euler's function.

Let $n=2m$. Then $\omega^2$ is a primitive $m$-th root of unity, and
\[
N_{K/\mathbb Q}(1+\omega)=\frac{N_{K/\mathbb Q}(1-\omega^2)}{N_{K/\mathbb Q}(1-\omega)}
=\frac{\Phi_m(1)^{[K:\mathbb Q(\omega^2)]}}{\Phi_n(1)} .
\]
If $m=\ell^e$ with $\ell$ an odd prime, then $\mathbb Q(\omega^2)=K$
because $\varphi(2m)=\varphi(m)$, and $n=2\ell^e$ has two prime factors,
so the norm is $\Phi_m(1)/\Phi_n(1)=\ell$. If $m=2^{e-1}$ with $e\ge2$,
then $[K:\mathbb Q(\omega^2)]=\varphi(2^e)/\varphi(2^{e-1})=2$ and the
norm is $2^2/2=2$. In both cases $1+\omega$ has norm different from
$\pm1$ and is not a unit.

The family is infinite: by Lemma~\ref{lem:compat} it contains every
$a/2^e$ with $a$ odd and $q_0\le a/2^e<1/3$, and every $a/(2\ell)$ with
$\ell$ an odd prime, $a$ odd, $\ell\nmid a$ and $q_0\le a/(2\ell)<1/3$.
For $q=1/10$ we have $n=2\cdot5$, and Lemma~\ref{lem:compat} gives
$p=1/5$.
\end{proof}

\begin{remark}
The same computation shows that $\xi$ is a unit for every other $n\ge3$:
for odd $n$ the elements $1-\omega^2$ and $1-\omega$ are conjugate, and for
$n=2m$ with $m$ divisible by two distinct primes both norms are $1$. So
for the other $n$ the number $2\cos(\pi g)$ may be an algebraic integer. For the remaining denominators the
angle $g$ can still be irrational; in fact $g$ is irrational for $226$ of
the $232$ compatible pairs with $n\le60$ outside the family
(Section~\ref{sec:comp}).
\end{remark}

For $q=1/10$ the coloring of Theorem~\ref{thm:main} is the following.

\begin{corollary}\label{cor:tenth}
For $q=1/10$, $p=1/5$ we have $r=(1+\sqrt5)/2$, $s=\sqrt{(5+\sqrt5)/10}$,
$\cos(\pi g)=s$, and $19/110<g<21/110$. The choice $k=11$, $L=110$
satisfies \eqref{eq:L}, and the coloring \eqref{eq:color} reads
\[
\begin{aligned}
C(A(x))&=(0,1,2,0,1,2,0,1,2,1)_{\lfloor110x\rfloor\bmod10},\\
C(B(x))&=(2,2,0,1,2,0,1,2,0,0)_{\lfloor110x+1/2\rfloor\bmod10}.
\end{aligned}
\]
No coloring $\lfloor3mx\rfloor\bmod3$ is proper.
\end{corollary}

\begin{proof}
Put $w=\sqrt5$ and $t=s^2=(5+w)/10$. The values of $r=R(1/10)$ and
$s=R(1/5)$ are the standard pentagon values, and $r^2+t-1=2rt$ by direct
expansion, so $\cos(\pi g)=(r^2+s^2-1)/(2rs)=s$. Let $\gamma=\pi g\in(0,\pi/2)$.
From $2<w<5/2$ we get $(3+w)/8<t<3/4$, that is $\cos^2(\pi/5)<\cos^2\gamma<\cos^2(\pi/6)$,
so $\pi/6<\gamma<\pi/5$. The Chebyshev polynomial $T_{11}$ is odd, and
evaluating it at $s$ in $\mathbb Q(w)(s)$ gives
\[
\cos(11\gamma)=s\Big(\frac{19}{25}+\frac{22w}{125}\Big)>0,\qquad
\cos^2(11\gamma)=\frac12+\frac{6469w}{31250}.
\]
Since $\cos^2(\pi/10)=(5+w)/8$,
\[
\cos^2(11\gamma)-\cos^2(\pi/10)=-\frac18+\frac{10251w}{125000}>\frac{4877}{125000}>0 .
\]
As $11\gamma\in(11\pi/6,11\pi/5)\subset(3\pi/2,5\pi/2)$ and $\cos(11\gamma)>\cos(\pi/10)>0$,
we get $|11\gamma-2\pi|<\pi/10$, which is $19/110<g<21/110$. Then
$Lg/2=55g\in(9.5,10.5)$, so $v=1$, $|\varepsilon|<1/2$ in \eqref{eq:L},
and $11\cdot1\equiv1\pmod{10}$. The two sequences are those of
Lemma~\ref{lem:palette} for $n=10$, where $m_-=2$ and $m_+=8$. The pair lies in the
family of Theorem~\ref{thm:family}, so $g$ is irrational, and $3\nmid10$
rules out every sector coloring by Proposition~\ref{prop:common}.
\end{proof}

\section{Computations}\label{sec:comp}

The following computations use exact rational arithmetic, SymPy for
algebraic numbers and mpmath at $60$ digits; the code is in the
repository named under Data and code availability.
\begin{enumerate}
\item There are $367$ pairs $q=a/n<1/3$ with $n\le60$ and
$p\equiv\pm2q$, and $284$ of them are compatible. For each compatible
pair the minimal polynomial of $2\cos(\pi g)$ is computed as a factor of
a resultant, and Kronecker's theorem~\cite{Kro57} decides whether $g$ is
rational: an algebraic integer whose conjugates are all real and lie in
$[-2,2]$ is $2\cos$ of a rational multiple of $\pi$. The number $g$ is
rational for exactly six pairs, $q=1/5,1/7,2/7,1/9,2/9,1/12$, with
$g=1/2,1/3,2/3,2/9,5/9,1/12$. Of the $284$ compatible pairs, $52$ lie in
the family of Theorem~\ref{thm:family}.
\item For each of the remaining $278$ pairs, the least $k$ allowed in
Section~\ref{sec:proof} is computed, and the coloring \eqref{eq:color}
is compared with the edges \eqref{eq:edges} arc by arc. For every arc and
every shift $\pm q$, $\pm p$, $\pm g/2$, the arcs met by the shifted arc
are found as intervals and their colors compared. No edge is
monochromatic. The largest $L$ is $2{,}445{,}927$, and the smallest
distance from an arc endpoint to a shifted arc endpoint is
$1.7\times10^{-3}$ of an arc. For $q=1/10$ the least $k$ is $11$. The
two colorings of Remark~\ref{rem:hyp} have $64$ and $88$ monochromatic
pairs of arcs.
\end{enumerate}

\section{Questions}\label{sec:questions}

The six rational values of $g$ found for $n\le60$ all have small
denominators. Rationality of $g$ is a linear relation with rational
coefficients among cosines of rational multiples of $\pi$, and relations
of this kind were classified by Conway and Jones~\cite{CJ76}.

\emph{Question 1.} Are $q=1/5,1/7,2/7,1/9,2/9,1/12$ the only compatible
$q$ with $p\equiv\pm2q$ for which $g$ is rational? A positive answer
would extend Theorem~\ref{thm:main} to every other such pair.

\emph{Question 2.} When $p\not\equiv\pm2q$ the finite graph $T_n$ is
replaced by another one. Which rational compatible pairs with $g$
irrational have Borel chromatic number $3$?

\emph{Question 3.} Bound the least $k$ allowed in
Section~\ref{sec:proof}, and hence the number of arcs, in terms of $n$
and $g$.

\section*{Acknowledgments and funding}

The research presented in this paper was supported by the European
Research Council (ERC) under the European Union's Horizon Europe research
and innovation programme (grant agreement No.~101041711) and by the Simons
Foundation as part of the Collaboration on the Mathematical and Scientific
Foundations of Deep Learning. It was also supported by Heights Labs, by
Convex Nexus Capital and by the Israel Science Foundation (ISF grants
2258/19 and 4101/25).

This work was done by multiple AI agents, Claude (Anthropic) and Codex
(OpenAI), connected to Math Brain~\cite{MathBrainRepo} and working
together inside the Brain framework~\cite{BrainFramework}. The method is
described in~\cite{MathBrainPaper}.

\section*{Declaration of competing interest}

The author holds equity in Heights Labs and in Convex Nexus Capital and
has a management role at Convex Nexus Capital.

\section*{Declaration of generative AI and AI-assisted technologies}

This paper is part of a broader research project on proving open
problems with AI systems, carried out with Claude (Anthropic) and Codex (OpenAI) connected
to Math Brain~\cite{MathBrainRepo} inside the Brain
framework~\cite{BrainFramework}, as described in~\cite{MathBrainPaper}.
Codex (model GPT-6.1-sol) found and proved Theorem~\ref{thm:main},
Lemma~\ref{lem:palette}, Proposition~\ref{prop:common} and
Corollary~\ref{cor:tenth}, and a separate Codex run audited them. Claude
(Opus~5.5) checked these proofs, found and proved Lemma~\ref{lem:compat},
Theorem~\ref{thm:family} and Remark~\ref{rem:radial}, wrote the code of
Section~\ref{sec:comp} and wrote the text. The rose window graphs are
Wilson's~\cite{Wil08}. Radial colorings of annuli in the unit-distance
graph were studied by Alm and Manske~\cite{AM14}, and the sector
colorings of Proposition~\ref{prop:common} are of the kind used by
Chybowska-Sok\'o\l, Junosza-Szaniawski and W\k{e}sek~\cite{CJW22}. The
facts on cyclotomic units behind Theorem~\ref{thm:family} are from
Washington's book~\cite{Was97}.
The author directed the work and is
responsible for the content of the paper and for decisions concerning its
dissemination.

\section*{Data and code availability}

The script \texttt{verify.py} and its output are available at
\url{https://github.com/yspennstate/unit-distance-concentric-circles}.

\begin{thebibliography}{99}

\bibitem{AM14}
J.~F. Alm and J.~Manske, \emph{On radial colorings of annuli},
Australas. J. Combin. \textbf{60} (2014), 270--278, arXiv:1305.5229.

\bibitem{CMT16}
C.~T. Conley, A.~S. Marks and R.~D. Tucker-Drob, \emph{Brooks's theorem
for measurable colorings}, Forum Math. Sigma \textbf{4} (2016), e16,
doi:10.1017/fms.2016.14.

\bibitem{CJ76}
J.~H. Conway and A.~J. Jones, \emph{Trigonometric diophantine equations
(On vanishing sums of roots of unity)}, Acta Arith. \textbf{30} (1976),
229--240, doi:10.4064/aa-30-3-229-240.

\bibitem{CJW22}
J.~Chybowska-Sok\'o\l, K.~Junosza-Szaniawski and K.~W\k{e}sek,
\emph{Coloring distance graphs on the plane}, Discrete Math.
\textbf{346} (2023), no.~7, 113441, doi:10.1016/j.disc.2023.113441;
arXiv:2201.04499 (the numbering cited is that of the arXiv version).

\bibitem{GR21}
J.~Greb\'ik and V.~Rozho\v{n}, \emph{Classification of local problems on
paths from the perspective of descriptive combinatorics}, preprint (2021),
arXiv:2103.14112.

\bibitem{KST99}
A.~S. Kechris, S.~Solecki and S.~Todorcevic, \emph{Borel chromatic
numbers}, Adv. Math. \textbf{141} (1999), 1--44,
doi:10.1006/aima.1998.1771.

\bibitem{Kro57}
L.~Kronecker, \emph{Zwei S\"atze \"uber Gleichungen mit ganzzahligen
Coefficienten}, J. Reine Angew. Math. \textbf{53} (1857), 173--175,
doi:10.1515/crll.1857.53.173.

\bibitem{Soi09}
A.~Soifer, \emph{The Mathematical Coloring Book}, Springer, New York,
2009, doi:10.1007/978-0-387-74642-5.

\bibitem{Was97}
L.~C. Washington, \emph{Introduction to Cyclotomic Fields}, 2nd ed.,
Graduate Texts in Mathematics \textbf{83}, Springer, New York, 1997.

\bibitem{Wei20}
F.~Weilacher, \emph{Descriptive chromatic numbers of locally finite and
everywhere two-ended graphs}, Groups Geom. Dyn. \textbf{16} (2022),
141--152, doi:10.4171/GGD/643; arXiv:2004.02316.

\bibitem{Wil08}
S.~Wilson, \emph{Rose window graphs}, Ars Math. Contemp. \textbf{1}
(2008), 7--19, doi:10.26493/1855-3974.13.5bb.

\bibitem{MathBrainRepo} Y.~Shmalo, \emph{Math Brain}, public repository, 2026, \url{https://github.com/yspennstate/math-brain-public/tree/8c9a6bd478c4ec4082347151b7ae16db50c917dd}.

\bibitem{MathBrainPaper} Y.~Shmalo, \emph{Math Brain: shared, auditable memory for multi-agent mathematical research}, October 2026, doi:10.5281/zenodo.24150837,
\url{https://github.com/yspennstate/math-brain-public/tree/main/papers/math-brain-2026-10-24}.

\bibitem{BrainFramework} Y.~Shmalo, \emph{Brain Framework}, public software repository, 2026, \url{https://github.com/yspennstate/brain-framework-public/tree/5629e644cdd510173ae97e9aa92642dabd22dc57}.

\end{thebibliography}

\end{document}
